% ============================================================================
% PART II: PHYSICAL APPLICATIONS
% Chapter 12: Experiments on Anyonic Braiding
% ============================================================================
%
% Purpose
%   Give each major experiment 2-3 pages: device, observable, theoretical
%   prediction tested, what was actually measured, caveats. Required honesty
%   boxes for non-abelian claims.
%
% Status: NEW WRITING. Guided by BSIEW tasks 12_1--12_4.
%
% Length target: ~10 pages
%
% Section plan
%   12.1  Nakamura 2020: nu=1/3 Fabry-Perot interferometer
%   12.2  Werkmeister 2024: graphene interferometer
%   12.3  Andersen 2023: non-abelian braiding on superconducting processor
%   12.4  Bartolomei 2020 and shot-noise experiments
%   12.5  Figures
%
% Owner: Helena
% Stage: V
%
% BSIEW task files:
%   12_1_nakamura_2020.md
%   12_2_werkmeister_2024.md
%   12_3_andersen_2023.md
%   12_4_experiment_figures.md
%
% References
%   Nakamura et al. 2020 (\cite{Nakamura2020AnyonicBraiding})
%   Werkmeister et al. 2024 (\cite{Werkmeister2024GrapheneInterferometer})
%   Andersen et al. 2023 (\cite{Andersen2023NonAbelianBraiding})
%   Bartolomei et al. 2020 (\cite{Bartolomei2020AnyonCollisions})
%   Saminadayar et al. 1997 (\cite{Saminadayar1997FractionalCharge})
%   de-Picciotto et al. 1997 (\cite{dePicciotto1997FractionalCharge})
%   Camino-Zhou-Goldman 2005 (\cite{CaminoZhouGoldman2005ABSuperperiod})
%   Venkatachalam et al. (\cite{VenkatachalamHartPfeifferWestYacobyGrapheneInterferometers})
%   Rosenow-Halperin (\cite{RosenowHalperinInterferometryPlaceholder})
%
% Dependencies
%   - Ch 11 (theoretical predictions cited back).
%   - Ch 10 (anyon data for comparison).
%
% Writing notes
%   - For each experiment: device -> observable -> prediction -> measurement.
%   - Andersen 2023 REQUIRES honesty box: synthetic emulation, not natural phase.
%   - Bartolomei 2020 is a new subsection not in the BSIEW plan; it expands
%     what was a parenthetical citation in the old Sec 5.2.

\section{Experiments on Anyonic Braiding}
\label{sec:experiments}

\subsection{Nakamura 2020: $\nu = 1/3$ Fabry--P\'erot Interferometer}
\label{subsec:nakamura-2020}

% TODO (Task 12_1): ~3 pages.
% 1. Device: GaAs/AlGaAs heterostructure, Fabry-Perot with two QPCs.
% 2. Observable: interference oscillation as magnetic flux varies; discrete
%    phase jump when bulk quasiparticle enters interference region.
% 3. Measured value: Delta phi ~ 2 pi/3, consistent with Laughlin braiding.
% 4. Theoretical prediction: from Ch 11 Sec 11.3 and Part I Wilson-loop:
%    e^{2 pi i / m} = e^{2 pi i/3} for nu = 1/3.
% 5. Statement: first direct observation of anyonic braiding (not inference).
% Cite: Nakamura et al. Nature Physics 16 (2020) 931.

\subsection{Werkmeister 2024: Graphene Interferometer}
\label{subsec:werkmeister-2024}

% TODO (Task 12_2): ~2 pages.
% 1. Platform: hBN-encapsulated graphene, cleaner edges, tunable filling.
% 2. Claims: detection of anyonic statistics with improved phase resolution;
%    telegraph-noise signature.
% 3. Novelty: different material platform corroborates anyonic interpretation;
%    improved noise floor.
% Cite: Werkmeister et al. 2024.

\subsection{Andersen 2023: Non-Abelian Braiding on a Superconducting Processor}
\label{subsec:andersen-2023}

% TODO (Task 12_3): ~2 pages + REQUIRED HONESTY BOX.
% 1. Platform: superconducting qubit array programmed to realize Z_2 toric
%    code + defects carrying non-abelian (Ising-like) braiding.
% 2. What was demonstrated: initialization, braiding, measurement of
%    non-abelian anyons (graph vertices with Ising fusion).
% 3. What was NOT demonstrated: not a natural topologically ordered phase.
%    Hamiltonian is programmed in software. Decoherence limits fidelity.
%
% REQUIRED HONESTY BOX (tcolorbox or visually distinguished):
% "This experiment realizes the braid-group algebra of non-abelian anyons on
%  a quantum processor, but it is not a discovery of non-abelian topological
%  order in a natural material. The distinction matters: natural topological
%  order would provide passive protection from decoherence, whereas the
%  synthetic realization relies on active error correction during the braiding
%  protocol."
%
% Cite: Andersen et al. 2023, Nature.

\subsection{Bartolomei 2020 and Shot-Noise Experiments}
\label{subsec:bartolomei-shot-noise}

% TODO (new subsection, not in BSIEW plan): ~1.5 pages.
% Currently only parenthetical citations in Ch 11. Expand:
% 1. Bartolomei 2020: collider-style noise measurement, what the
%    observable is, how noise evidence differs from interferometry.
% 2. Earlier shot-noise: Saminadayar 1997, de-Picciotto 1997.
%    Fractional charge e/3 from current fluctuations.
% 3. How shot-noise complements interferometry: charge vs. statistics.
% Cite: Bartolomei et al. 2020, Saminadayar 1997, de-Picciotto 1997.

\subsection{Figures}
\label{subsec:experiment-figures}

% TODO (Task 12_4): 3 figures.
% 1. Interferometer schematic for Nakamura 2020.
% 2. Graphene device for Werkmeister 2024.
% 3. Qubit-array diagram for Andersen 2023.
%
% Approach: obtain permission or redraw schematically.
% Each caption cites original with DOI.
% If redrawn: "schematic based on [citation]."
